\documentclass[11pt]{article}

\usepackage[margin=1in]{geometry}
\usepackage{amsmath,amssymb}

\title{A Scattering-Candidate Reading of the Local Value Scale}
\author{Working note in \texttt{grh/}}
\date{2026-04-23}

\begin{document}

\maketitle

\begin{abstract}
This note records a candidate-object refinement suggested by the current RH
draft. The positive scalar \(S(m)\) used in the local value channel is more
naturally read as a positive strip-edge zero kernel than as a surrogate for a
raw critical-line value. This suggests that future Dirichlet and Ramanujan-
\(\tau\) generalizations should target a completed-object scattering package,
not merely a real critical-line observable. The note does not claim that the
draft already proves the full completed-quotient identity for \(q\), nor that
any non-zeta realization is currently available.
\end{abstract}

\section{The kernel}

The draft uses the positive kernel
\[
f_{\beta,\gamma}(m)=
\frac{1-\beta}{(1-\beta)^2+(2m-\gamma)^2}
+
\frac{\beta}{\beta^2+(2m-\gamma)^2}.
\]

For a zero \(\rho=\beta+i\gamma\), this can be rewritten as
\[
f_{\beta,\gamma}(m)
=
\Re\left(\frac{1}{1+2im-\rho}-\frac{1}{2im-\rho}\right).
\]

So the manuscript scalar
\[
S(m)=\sum_\rho f_{\beta_\rho,\gamma_\rho}(m)
\]
is a positive strip-edge zero contribution.

\section{Interpretation}

This kernel strongly suggests a completed-object logarithmic-derivative
difference, heuristically of strip-edge type,
\[
\Re\sum_\rho\left(\frac{1}{1+2im-\rho}-\frac{1}{2im-\rho}\right).
\]

Equivalently, one is led toward a scattering-style phase derivative attached to
a completed object, rather than toward a raw critical-line value or a mere
reality normalization.

The honest boundary is important: this note does \emph{not} claim that the
current manuscript already proves the exact formula-level identity
\[
q(t)=\frac12\Im\left(\frac{d}{dt}\log \mathcal Q(t)\right)
\]
for some explicitly named completed quotient \(\mathcal Q(t)\). The present
point is narrower. The zero kernel behind \(S(m)\) has the right strip-edge
shape to make such a reading natural.

\section{Archive provenance for the quotient candidate}

The repository's chat archive contains a precise candidate source formula for
this interpretation, although it is not yet promoted into the canonical RH
draft. In

\begin{quote}
\texttt{paper/chats/20260410-043629-69d87e03-a5c8-83a5-9f21-1062e8b6d064-riemann-hypothesis-insight.md:5528-5598}
\end{quote}

the proposed quotient is
\[
\phi(s)=\frac{\Lambda(2s-1)}{\Lambda(2s)},
\qquad
\phi\!\left(\tfrac12+it\right)=e^{-2i\Phi(t)}.
\]

In that archived derivation, a zero \(\rho=\beta+i\gamma\) contributes a pole at
\(\rho/2\) and a zero at \((1+\rho)/2\), which produces exactly the strip-edge
kernel used in the present note.

This is valuable provenance for future work in \texttt{grh/}, but it should still be
treated as a candidate source theorem rather than as an already-canonical part
of \texttt{paper/proof_of_rh.tex}.

\section{Why this matters for GRH work}

The earlier GRH notes treated candidate objects mainly as channel-choice
problems: real critical-line scalar, paired scalar, or matrix package. The
kernel identity sharpens that discussion.

What the non-zeta program should try to reproduce is not merely a real
observable. It should try to reproduce a background-subtracted scalar with the
same structural role as \(S(m)\):

\begin{enumerate}
\item real;
\item nonnegative;
\item built from strip-edge zero data of a completed object;
\item serving as the coefficient of the leading local value channel.
\end{enumerate}

That is a more faithful target than a raw rotated completed value.

\section{Dirichlet and \texorpdfstring{$\tau$}{tau}}

For primitive Dirichlet \(L\)-functions, this viewpoint sharpens the earlier
channel discussion. The real question is not only which object is real on the
critical line. The sharper question is which completed-object package can supply
a positive scalar analogue of \(S(m)\) inside the local deformation.

For Ramanujan-\(\tau\), the same viewpoint says that self-duality helps only at
the level of candidate-object selection. The actual localization package is
still missing.

\section{What remains missing}

The scattering reading does not by itself provide any of the hard realization
work. In particular, it does not yet supply:

\begin{enumerate}
\item an explicit completed-quotient theorem for the manuscript's \(q(t)\);
\item a Dirichlet realization of the corrected local block;
\item a Ramanujan-\(\tau\) realization of the corrected local block;
\item microscopic denominator comparability in those families;
\item preserved whitening hierarchy and boundary control in those families.
\end{enumerate}

So the gain here is conceptual but still real: better candidate-object
selection.

\section{Use}

This note should be read together with the other \texttt{grh/} notes as follows:

\begin{enumerate}
\item \texttt{local_odd_projector_note.tex} isolates the unconditional odd-germ layer;
\item \texttt{portability_note.tex} records the broader conservative portability view;
\item the present note sharpens the correct candidate object for the missing
realization theorem.
\end{enumerate}

\end{document}
